%==============================================================================
% SUMMARY
%
% Pauta: "Traducción al idioma inglés del Resumen."
%
% En la pauta aparece rotulado "1. Summary", el mismo número que el Resumen
% (es una errata del documento original: el capítulo siguiente es "2.
% Introducción"). Por eso aquí va como capítulo SIN número, listado en el
% Índice General justo después del Resumen.
%
% Si tu comisión prefiere que lleve número propio, reemplaza la línea
%   \capitulosinnumero{Summary}
% por
%   \chapter{Summary}
% y recuerda que la Introducción pasará a ser el capítulo 3.
%==============================================================================

\capitulosinnumero{Summary}

\begin{otherlanguage}{english}

Write here the English translation of the Resumen. It must mirror the Spanish
abstract: objectives, materials and methods, results, discussion and
conclusions, within a single letter-size page.

\vspace{1em}
\noindent\textbf{Keywords:} keyword one, keyword two, keyword three,
keyword four, keyword five.

\end{otherlanguage}
